\documentclass[12pt]{article}
\usepackage[utf8]{inputenc}
\usepackage[T1]{fontenc}
\usepackage[margin=1in]{geometry}
\usepackage{setspace}
\usepackage{xcolor}
\usepackage[colorlinks=true, urlcolor=blue!60!black, linkcolor=blue!60!black]{hyperref}
\setlength{\parindent}{2em}
\setlength{\parskip}{1.25em}
\renewcommand{\baselinestretch}{1.0}
\begin{document}
\begin{center}
{\LARGE \textit{Lawrence v. Texas}, 539 U.S. 558 (2003)\par}
\vspace{0.8em}
{\large Case Brief\par}
\end{center}
\vspace{1em}

\section{Memory Jogger}
In \href{https://www.courtlistener.com/opinion/130160/lawrence-v-texas/}{\textit{Lawrence v. Texas}}, 539 U.S. 558 (2003), the Supreme Court overruled \textit{Bowers v. Hardwick} and held 6--3 that a Texas law criminalizing private, consensual same-sex sexual intimacy violates the Fourteenth Amendment's Due Process Clause, rejecting majoritarian moral disapproval as a legitimate state interest and deepening the methodological divide over substantive due process.

\section{Detailed Case Facts}
On September 17, 1998, Harris County Sheriff's deputies in Houston, Texas were dispatched to the apartment of petitioner John Geddes Lawrence in response to a reported weapons disturbance (\href{https://www.law.cornell.edu/supremecourt/text/02-102}{539 U.S. at 562--63}; \href{https://www.tshaonline.org/handbook/entries/lawrence-v-texas}{Texas State Historical Association}). The report had been falsely phoned in by an acquaintance, Robert Eubanks, who later pleaded no contest to filing a false police report (\href{https://www.law.cornell.edu/wex/lawrence_v._texas}{Cornell LII Wex}). Upon entering Lawrence's apartment---an entry whose Fourth Amendment legality was not challenged before the Supreme Court---officers observed Lawrence and another adult man, petitioner Tyron Garner, engaging in a private, consensual sexual act (\href{https://www.law.cornell.edu/supremecourt/text/02-102}{539 U.S. at 563}).

The officers arrested both men, held them in custody overnight, and charged them with violating Texas's ``Homosexual Conduct'' statute, Tex. Penal Code Ann. \S~21.06(a) (2003), enacted in 1973 when Texas revised its criminal code to decriminalize opposite-sex sodomy while retaining criminal penalties for same-sex couples (\href{https://www.law.cornell.edu/supremecourt/text/02-102}{539 U.S. at 563, 570}). Section 21.06(a) made it a Class C misdemeanor if a person ``engages in deviate sexual intercourse with another individual of the same sex,'' with \S~21.01(1) defining ``deviate sexual intercourse'' as oral or anal genital contact or penetration of the genitals or anus with an object (\href{https://www.law.cornell.edu/supremecourt/text/02-102}{539 U.S. at 563}). The criminal complaints described petitioners' offense as ``deviate sexual intercourse, namely anal sex, with a member of the same sex (man)'' (\href{https://www.law.cornell.edu/supremecourt/text/02-102}{539 U.S. at 563}). Although Texas rarely enforced \S~21.06(a) against consenting adults in private, a conviction carried a permanent criminal record, potential occupational licensing restrictions in Texas, and mandatory sex-offender registration in at least four other states (\href{https://www.law.cornell.edu/supremecourt/text/02-102}{539 U.S. at 575--76, 581}).

\section{Procedural History}
Lawrence and Garner were initially convicted before a Harris County Justice of the Peace after pleading no contest (\href{https://www.law.cornell.edu/wex/lawrence_v._texas}{Cornell LII Wex}). Exercising their statutory right to a trial \textit{de novo} in the Harris County Criminal Court, they moved to quash the complaints on the ground that \S~21.06(a) violated the Equal Protection Clause of the Fourteenth Amendment and the Equal Rights Amendment of the Texas Constitution, Tex. Const., Art. 1, \S~3a (\href{https://www.law.cornell.edu/supremecourt/text/02-102}{539 U.S. at 563}). The county criminal court rejected those constitutional challenges; petitioners then entered pleas of \textit{nolo contendere}, and each was fined \$200 and assessed court costs of \$141.25 (\href{https://www.law.cornell.edu/supremecourt/text/02-102}{539 U.S. at 563}).

On appeal, a three-judge panel of the Texas Court of Appeals for the Fourteenth District initially reversed 2--1 on June 8, 2000, holding that \S~21.06(a) violated the Texas Equal Rights Amendment (\href{https://www.law.cornell.edu/wex/lawrence_v._texas}{Cornell LII Wex}). The Fourteenth Court of Appeals granted rehearing \textit{en banc}, however, and on March 15, 2001 affirmed the convictions by a 7--2 vote, rejecting petitioners' federal Equal Protection and Due Process challenges and treating \textit{Bowers v. Hardwick}, 478 U.S. 186 (1986), as controlling on substantive due process (\textit{Lawrence v. State}, 41 S.W.3d 349 (Tex. App. 2001); \href{https://www.law.cornell.edu/supremecourt/text/02-102}{539 U.S. at 563}). The Texas Court of Criminal Appeals denied discretionary review on April 17, 2002 (\href{https://www.law.cornell.edu/wex/lawrence_v._texas}{Cornell LII Wex}). The U.S. Supreme Court granted certiorari on December 2, 2002 (537 U.S. 1044) on three questions: (1) whether \S~21.06(a) violated Equal Protection; (2) whether petitioners' convictions violated Due Process; and (3) whether \textit{Bowers v. Hardwick} should be overruled (\href{https://www.oyez.org/cases/2002/02-102}{Oyez}). The case was argued on March 26, 2003 and decided on June 26, 2003 (\href{https://www.oyez.org/cases/2002/02-102}{Oyez}).

\section{Judicial Votes}
The Supreme Court reversed and remanded by a 6--3 vote (\href{https://www.law.cornell.edu/supremecourt/text/02-102}{539 U.S. at 561, 579}):
\begin{itemize}
\item \textbf{Majority (5 Justices):} Justice Anthony M. Kennedy delivered the opinion of the Court, joined by Justices John Paul Stevens, David H. Souter, Ruth Bader Ginsburg, and Stephen G. Breyer.
\item \textbf{Concurrence in the Judgment (1 Justice):} Justice Sandra Day O'Connor filed an opinion concurring in the judgment, invalidating the Texas law under the Equal Protection Clause without overruling \textit{Bowers}.
\item \textbf{Dissents (3 Justices):} Justice Antonin Scalia filed a dissenting opinion, joined by Chief Justice William H. Rehnquist and Justice Clarence Thomas. Justice Clarence Thomas also filed a brief separate dissenting opinion.
\end{itemize}

\section{Holding}
\textbf{Judgment:} The judgment of the Court of Appeals for the Texas Fourteenth District was reversed, and the case was remanded for further proceedings not inconsistent with the Court's opinion (\href{https://www.law.cornell.edu/supremecourt/text/02-102}{539 U.S. at 579}).

\textbf{Rule Announced:} The Fourteenth Amendment's Due Process Clause protects the liberty of consenting adults to engage in private, noncommercial intimate sexual conduct in the home without criminal punishment; majoritarian moral disapproval alone is not a legitimate state interest sufficient to justify criminalizing such conduct, and \textit{Bowers v. Hardwick}, 478 U.S. 186 (1986), is overruled (\href{https://www.courtlistener.com/opinion/130160/lawrence-v-texas/}{539 U.S. at 578}).

\textbf{Holding vs. Dicta:} The binding holding invalidates criminal prohibitions on private, consensual adult sexual intimacy under the Due Process Clause and expressly overrules \textit{Bowers}. By contrast, the Court's closing caveats noting that the case did not involve ``minors,'' ``persons who might be injured or coerced,'' ``public conduct or prostitution,'' or ``whether the government must give formal recognition to any relationship that homosexual persons seek to enter'' are limiting observations (dicta as to unpresented issues) rather than holdings authorizing or forbidding restrictions in those domains (\href{https://www.law.cornell.edu/supremecourt/text/02-102}{539 U.S. at 578}). Similarly, the majority's discussion of foreign and international authorities---including the British Wolfenden Report and the European Court of Human Rights decision in \textit{Dudgeon v. United Kingdom}---served as persuasive context rebutting \textit{Bowers}'s appeal to Western civilization rather than binding constitutional authority (\href{https://www.law.cornell.edu/supremecourt/text/02-102}{539 U.S. at 572--73, 576--77}).

\section{Analysis of Opinions}

\subsection{Majority Opinion (Justice Kennedy)}
Justice Kennedy began by framing liberty in both ``spatial and more transcendent dimensions,'' emphasizing that the Constitution protects an ``autonomy of self that includes freedom of thought, belief, expression, and certain intimate conduct'' (\href{https://www.law.cornell.edu/supremecourt/text/02-102}{539 U.S. at 562}). Tracing substantive due process doctrine from \textit{Pierce v. Society of Sisters} and \textit{Meyer v. Nebraska} through \textit{Griswold v. Connecticut}, \textit{Eisenstadt v. Baird}, \textit{Roe v. Wade}, and \textit{Carey v. Population Services International}, Kennedy concluded that \textit{Bowers v. Hardwick} had fatally misframed the liberty interest at a hyper-specific level of generality by asking whether there was a ``fundamental right upon homosexuals to engage in sodomy'' (\href{https://www.law.cornell.edu/supremecourt/text/02-102}{539 U.S. at 564--67}). Reducing private intimacy to a single sexual act demeaned the personal relationship at stake just as reducing marriage to sexual intercourse would demean married couples (\href{https://www.law.cornell.edu/supremecourt/text/02-102}{539 U.S. at 567}).

Turning to the substantive due process methodology war, Kennedy dismantled \textit{Bowers}'s historical premises and refused to apply the rigid backward-looking test of \textit{Washington v. Glucksberg}, 521 U.S. 702 (1997). He observed that early American sodomy statutes targeted nonprocreative sex generally rather than homosexuals as a distinct class, were rarely enforced against consenting adults in private, and did not single out same-sex couples until the 1970s (\href{https://www.law.cornell.edu/supremecourt/text/02-102}{539 U.S. at 568--70}). Crucially, Kennedy quoted his own concurrence in \textit{County of Sacramento v. Lewis} that ``history and tradition are the starting point but not in all cases the ending point of the substantive due process inquiry,'' pointing instead to an ``emerging awareness'' over the preceding half-century---reflected in the 1955 Model Penal Code, state repeals reducing sodomy bans from 25 states in 1986 to 13 in 2003, the Wolfenden Report, and \textit{Dudgeon v. United Kingdom} (\href{https://www.law.cornell.edu/supremecourt/text/02-102}{539 U.S. at 571--73}).

Kennedy declined to decide the case solely on Equal Protection grounds because leaving \textit{Bowers} intact would allow states to criminalize both same-sex and opposite-sex sodomy and would perpetuate the stigma and ``invitation to subject homosexual persons to discrimination'' (\href{https://www.law.cornell.edu/supremecourt/text/02-102}{539 U.S. at 574--75}). Invoking \textit{Planned Parenthood of Southeastern Pa. v. Casey} and \textit{Romer v. Evans} as eroding \textit{Bowers}'s foundations, and finding no detrimental reliance under \textit{stare decisis}, Kennedy adopted Justice Stevens's \textit{Bowers} dissent: traditional moral disapproval is not a sufficient basis to uphold a criminal prohibition on private intimate choices (\href{https://www.law.cornell.edu/supremecourt/text/02-102}{539 U.S. at 576--78}). Notably, Kennedy avoided labeling the right ``fundamental'' or explicitly announcing strict scrutiny, instead declaring that the Texas statute ``furthers no legitimate state interest which can justify its intrusion into the personal and private life of the individual'' (\href{https://www.law.cornell.edu/supremecourt/text/02-102}{539 U.S. at 578}).

\subsection{Concurrence in the Judgment (Justice O'Connor)}
Justice O'Connor, who had joined the \textit{Bowers} majority, refused to overrule \textit{Bowers} or rely on substantive due process (\href{https://www.law.cornell.edu/supremecourt/text/02-102}{539 U.S. at 579}). Instead, she concluded that \S~21.06(a) violated the Fourteenth Amendment's Equal Protection Clause under ``a more searching form of rational basis review'' drawn from \textit{Department of Agriculture v. Moreno}, \textit{City of Cleburne v. Cleburne Living Center}, and \textit{Romer v. Evans} (\href{https://www.law.cornell.edu/supremecourt/text/02-102}{539 U.S. at 580}). Because Texas punished only same-sex sodomy while leaving identical opposite-sex conduct legal, the statute targeted gay persons as a class, branded them as criminals, and sanctioned collateral discrimination (\href{https://www.law.cornell.edu/supremecourt/text/02-102}{539 U.S. at 581--83}). Moral disapproval of a politically unpopular group, like a bare desire to harm that group, is not a legitimate governmental interest under Equal Protection (\href{https://www.law.cornell.edu/supremecourt/text/02-102}{539 U.S. at 582--83}). O'Connor added that a facially neutral sodomy ban would likely not survive democratic politics and suggested that other distinctions, such as preserving ``the traditional institution of marriage,'' could rest on legitimate interests beyond mere moral disapproval (\href{https://www.law.cornell.edu/supremecourt/text/02-102}{539 U.S. at 584--85}).

\subsection{Principal Dissent (Justice Scalia)}
Justice Scalia, joined by Chief Justice Rehnquist and Justice Thomas, launched a blistering attack on the majority's treatment of \textit{stare decisis}, substantive due process methodology, and judicial restraint. First, he accused the Court of applying a manipulative double standard by abandoning \textit{Casey}'s heightened \textit{stare decisis} deference to ``intensely divisive'' precedents when overruling \textit{Bowers}, while refusing to apply the same erosion-and-criticism criteria to \textit{Roe v. Wade} (\href{https://www.law.cornell.edu/supremecourt/text/02-102}{539 U.S. at 586--91}). He derided \textit{Casey}'s ``sweet-mystery-of-life passage'' as ``the passage that ate the rule of law'' (\href{https://www.law.cornell.edu/supremecourt/text/02-102}{539 U.S. at 588}).

Second, Scalia emphasized that the majority never declared homosexual sodomy a ``fundamental right'' under \textit{Washington v. Glucksberg}---because it could not show such a right was ``deeply rooted in this Nation's history and tradition''---and instead applied ``an unheard-of form of rational-basis review'' (\href{https://www.law.cornell.edu/supremecourt/text/02-102}{539 U.S. at 586, 593--94}). He argued that an ``emerging awareness'' is by definition not deeply rooted, criticized reliance on foreign law as ``dangerous dicta,'' and warned that holding majoritarian morality illegitimate ``effectively decrees the end of all morals legislation,'' calling into question laws against bigamy, adult incest, prostitution, adultery, obscenity, and same-sex marriage bans (\href{https://www.law.cornell.edu/supremecourt/text/02-102}{539 U.S. at 590, 598--99, 604--05}). Rejecting O'Connor's Equal Protection rationale as well, Scalia contended that \S~21.06(a) applied equally to men and women and that the Court had ``taken sides in the culture war'' rather than leaving moral change to the democratic process (\href{https://www.law.cornell.edu/supremecourt/text/02-102}{539 U.S. at 599--603}).

\subsection{Separate Dissent (Justice Thomas)}
Justice Thomas joined Justice Scalia's dissent in full and wrote two paragraphs quoting Justice Potter Stewart's dissent in \textit{Griswold v. Connecticut} to call the Texas law ``uncommonly silly'' (\href{https://www.law.cornell.edu/supremecourt/text/02-102}{539 U.S. at 605}). While stating that he would vote to repeal \S~21.06(a) as a Texas legislator, Thomas concluded that he could find neither a general ``right of privacy'' nor a ``liberty of the person both in its spatial and more transcendent dimensions'' anywhere in the Constitution, and therefore was not empowered to invalidate the law (\href{https://www.law.cornell.edu/supremecourt/text/02-102}{539 U.S. at 605--06}).

\section{Examples: Future Applications}
\begin{enumerate}
\item \textbf{Same-Side (Invalidated): State Criminal Ban on Private Consensual Fornication or Cohabitation.} Suppose a state enforces a statute criminalizing private, consensual sexual intercourse between unmarried heterosexual adults in their home, justified solely by traditional religious morality. This falls squarely on the same side as \textit{Lawrence}: it intrudes into private, noncommercial adult intimacy with full mutual consent and rests on moral disapproval rather than prevention of third-party harm.
\item \textbf{Same-Side (Invalidated): Criminal Prohibition on the Private Possession or Sale of Adult Sexual Devices.} Suppose a state prohibits the retail distribution of sexual devices used solely by consenting adults in private, asserting only an interest in ``safeguarding public morality.'' Following \textit{Lawrence}'s rejection of bare moral disapproval as a legitimate state interest intruding on private adult sexual autonomy (as later seen in post-\textit{Lawrence} Fifth Circuit litigation), the statute should fail because it burdens the same private sphere of consensual adult intimacy without preventing tangible harm.
\item \textbf{Opposite-Side (Upheld): Criminal Prosecution for Commercial Prostitution.} Suppose a state prosecutes an adult for exchanging sexual acts for money in a private hotel room. Even though the participants are consenting adults behind closed doors, this falls on the opposite side of \textit{Lawrence}: Justice Kennedy expressly distinguished ``prostitution'' and commercial conduct, and the state asserts concrete non-moral interests---preventing human trafficking, coercion, exploitation, and secondary criminal markets---that easily survive review.
\item \textbf{Opposite-Side (Upheld): Statutory Rape or Sexual Conduct Involving Minors or Coercive Relationships.} Suppose an adult challenges a state statute criminalizing sexual contact with a 15-year-old minor, or a prison guard challenges a statute banning sexual relations with an incarcerated inmate. This falls squarely on the opposite side: \textit{Lawrence} explicitly excluded ``minors'' and persons ``situated in relationships where consent might not easily be refused,'' where the state's compelling interest in protecting vulnerable individuals from exploitation overrides any liberty claim.
\item \textbf{Fence-Sitter: State Criminal Ban on Consensual Adult Incest or Polygamous Cohabitation.} Suppose two biological adult siblings, or three consenting adults in a private polyamorous household, challenge a state criminal statute prohibiting their private sexual or domestic arrangement. This is a classic fence-sitter: challengers can invoke \textit{Lawrence}'s broad autonomy principle and rejection of bare moral disapproval when no third party is directly harmed, while the state can argue that unlike same-sex sodomy, incest and polygamy implicate genetic risks, family-structure coercion concerns, and deeply entrenched institutional protections (``abuse of an institution the law protects'') that \textit{Lawrence} left open---and post-\textit{Dobbs} courts applying \textit{Glucksberg}'s history-and-tradition test would be far less receptive to extending \textit{Lawrence}.
\end{enumerate}

\section{Critique}
\textit{Lawrence v. Texas} stands at the epicenter of the Supreme Court's substantive due process ``methodology war,'' drawing sharp debate from originalist, libertarian, and progressive scholars alike.

From an \textbf{originalist and textualist perspective}, critics argue that Justice Kennedy abandoned the disciplined historical inquiry of \textit{Washington v. Glucksberg} without offering a coherent doctrinal replacement. Scholars Nelson Lund and John O. McGinnis condemned \textit{Lawrence} as ``judicial hubris,'' contending that the majority replaced constitutional text and history with an amorphous, unenumerated autonomy principle unconstrained by either traditional strict scrutiny or genuine rational-basis review (\href{https://repository.law.umich.edu/cgi/viewcontent.cgi?article=1719&context=mlr}{Lund \& McGinnis, 102 Mich. L. Rev. 1555 (2004)}). Similarly, Lino A. Graglia criticized the Court for reading Millian libertarianism into the Fourteenth Amendment and displacing representative self-government (\href{https://kb.osu.edu/bitstreams/f95c9b0e-cc7a-5c92-b702-7bf6e345f0a9/download}{Graglia, 65 Ohio St. L.J. 1141 (2004)}). This originalist critique later resurfaced in \href{https://supreme.justia.com/cases/federal/us/597/19-1392/}{\textit{Dobbs v. Jackson Women's Health Organization}}, 597 U.S. 215 (2022), where the Court restored \textit{Glucksberg}'s history-and-tradition test for abortion while distinguishing \textit{Lawrence}, and Justice Thomas urged the Court in concurrence to reconsider all substantive due process precedents, including \textit{Griswold}, \textit{Lawrence}, and \textit{Obergefell}.

From a \textbf{libertarian originalist perspective}, however, Randy E. Barnett praised \textit{Lawrence} as a ``libertarian revolution'' that properly shifted away from the rigid two-tier fundamental-rights framework toward a constitutional ``presumption of liberty,'' requiring the government to justify restrictions on harmless private conduct with a legitimate, non-moralistic purpose (\href{https://scholarship.law.georgetown.edu/cgi/viewcontent.cgi?article=1854&context=facpub}{Barnett, 2002--2003 Cato Sup. Ct. Rev. 21}).

From a \textbf{progressive and doctrinal perspective}, scholars celebrated \textit{Lawrence}'s vindication of LGBTQ dignity and its paving of the doctrinal path to \textit{United States v. Windsor} (2013) and \textit{Obergefell v. Hodges} (2015), yet also noted its doctrinal opacity. Laurence H. Tribe famously described \textit{Lawrence} as ``the `fundamental right' that dare not speak its name,'' praising how Kennedy wove Due Process and Equal Protection into an interlocking ``double helix'' of equal dignity even while eschewing formal tiers of scrutiny (\href{https://harvardlawreview.org/forum/vol-133/criminal-justice-user-fees-and-the-procedural-aspect-of-equal-justice/}{citing Tribe, 117 Harv. L. Rev. 1893 (2004)}). Cass Sunstein likewise highlighted the opinion's ambiguity---leaving lower courts uncertain whether \textit{Lawrence} rested on fundamental-rights autonomy, heightened rational basis, or constitutional desuetude (\href{https://scholarship.law.wm.edu/cgi/viewcontent.cgi?article=1647&context=wmborj}{citing Sunstein, 2003 Sup. Ct. Rev. 27}).

At bottom, \textit{Lawrence} exposes a foundational clash of \textbf{competing constitutional values}: individual autonomy, personal dignity, and anti-subordination against democratic self-governance, federalism, and the traditional police power of states to legislate public morality; and an evolving, principled reading of constitutional ``liberty'' against a fixed, history-and-tradition constraint designed to cabin judicial discretion.

\section{Key Quotations}
\begin{quote}
``Liberty protects the person from unwarranted government intrusions into a dwelling or other private places. In our tradition the State is not omnipresent in the home. And there are other spheres of our lives and existence, outside the home, where the State should not be a dominant presence. Freedom extends beyond spatial bounds. Liberty presumes an autonomy of self that includes freedom of thought, belief, expression, and certain intimate conduct. The instant case involves liberty of the person both in its spatial and more transcendent dimensions.''
\par\noindent---Justice Anthony M. Kennedy, Opinion of the Court, \href{https://www.law.cornell.edu/supremecourt/text/02-102}{539 U.S. at 562}.
\end{quote}

\begin{quote}
``Had those who drew and ratified the Due Process Clauses of the Fifth Amendment or the Fourteenth Amendment known the components of liberty in its manifold possibilities, they might have been more specific. They did not presume to have this insight. They knew times can blind us to certain truths and later generations can see that laws once thought necessary and proper in fact serve only to oppress. As the Constitution endures, persons in every generation can invoke its principles in their own search for greater freedom.''
\par\noindent---Justice Anthony M. Kennedy, Opinion of the Court, \href{https://www.law.cornell.edu/supremecourt/text/02-102}{539 U.S. at 578--79}.
\end{quote}

\begin{quote}
``Moral disapproval of this group, like a bare desire to harm the group, is an interest that is insufficient to satisfy rational basis review under the Equal Protection Clause.''
\par\noindent---Justice Sandra Day O'Connor, Concurring in the Judgment, \href{https://www.law.cornell.edu/supremecourt/text/02-102}{539 U.S. at 582}.
\end{quote}

\begin{quote}
``Today's opinion is the product of a Court, which is the product of a law-profession culture, that has largely signed on to the so-called homosexual agenda, by which I mean the agenda promoted by some homosexual activists directed at eliminating the moral opprobrium that has traditionally attached to homosexual conduct.''
\par\noindent---Justice Antonin Scalia, Dissenting, \href{https://www.law.cornell.edu/supremecourt/text/02-102}{539 U.S. at 602}.
\end{quote}

\begin{quote}
``I write separately to note that the law before the Court today `is \ldots{} uncommonly silly.' \textit{Griswold v. Connecticut}, 381 U.S. 479, 527 (1965) (Stewart, J., dissenting). If I were a member of the Texas Legislature, I would vote to repeal it.''
\par\noindent---Justice Clarence Thomas, Dissenting, \href{https://www.law.cornell.edu/supremecourt/text/02-102}{539 U.S. at 605}.
\end{quote}

\end{document}
